\chapter[VaR e CVaR]{VaR e CVaR: misurare e ottimizzare il rischio}
\label{cap:varcvar}

\noindent\textbf{Classe:} LP a scenari \hfill
\textbf{Script:} \texttt{python/lab13\_var\_cvar.py}

\section{Motivazione gestionale}

Molti modelli ottimizzano un valore medio. Scelta semplice, ma che può nascondere
perdite rare e molto elevate. VaR e CVaR descrivono la \emph{coda} della
distribuzione delle perdite e trasformano l'avversione al rischio in decisioni
quantitative. Il VaR risponde a ``qual è una soglia di perdita elevata?'';
il CVaR risponde anche a ``quanto perdiamo \emph{in media} quando quella soglia viene
superata?''.



Dopo le crisi finanziarie, ``quanto perdiamo negli scenari peggiori?'' è diventata
una domanda regolamentare oltre che gestionale. VaR e CVaR sono il linguaggio
standard con cui banche, assicurazioni e, sempre più, supply chain manager parlano
di rischio di coda. La sorpresa didattica è che il CVaR --- concetto
apparentemente sofisticato --- si ottimizza con un semplice LP.

\textbf{In questo capitolo impareremo a:} (1) definire e calcolare VaR e CVaR;
(2) usare la formulazione lineare di Rockafellar--Uryasev; (3) costruire portafogli
e piani logistici protetti dagli scenari peggiori; (4) riconoscere i limiti
statistici delle stime di coda.

\section{Definizioni e proprietà}

\paragraph{Il problema a parole.} Nei modelli di questo capitolo \emph{decidiamo}
come prima (quote di portafoglio, quantità, capacità prenotate), ma
\emph{l'obiettivo} cambia: non minimizzare la perdita \emph{media}, bensì una misura
della sua \emph{coda} --- quanto si perde negli scenari peggiori. \emph{I vincoli}
restano quelli del problema sottostante, più i vincoli lineari ausiliari che
definiscono il CVaR.

Sia $L$ la perdita aleatoria (valori alti $=$ esiti peggiori) e $\alpha \in (0,1)$
un livello di confidenza (tipicamente $0{,}90$--$0{,}99$).

\begin{modello}[VaR e CVaR]
\[
\var_\alpha(L) = \inf\{\eta \in \R : \Prob(L \le \eta) \ge \alpha\}
\qquad\quad
\cvar_\alpha(L) = \text{perdita media nella coda peggiore di massa } 1 - \alpha .
\]
\end{modello}

\begin{center}\small
\begin{tabular}{lp{10.2cm}}
\toprule
Convessità & il CVaR è convesso nelle decisioni (se $L$ lo è); il VaR in generale no \\
Subadditività & il CVaR premia la diversificazione; il VaR può violarla \\
Coda & il VaR non distingue due code con lo stesso quantile ma gravità diversa \\
Ottimizzazione & il CVaR ammette una formulazione LP a scenari; il VaR no \\
\bottomrule
\end{tabular}
\end{center}

\begin{modello}[Formulazione lineare di Rockafellar--Uryasev]
Con un insieme finito di scenari $S$, ciascuno con probabilità $\pi_s \ge 0$
($\sum_{s \in S} \pi_s = 1$) e perdita $\ell_s(\vet x)$ lineare nel vettore delle
decisioni $\vet x$:
\begin{align}
\min_{\vet x,\,\eta,\,\vet\xi}\;\; & \eta + \frac{1}{1-\alpha} \sum_{s \in S} \pi_s\, \xi_s \notag\\
\text{s.t.}\;\; & \xi_s \ge \ell_s(\vet x) - \eta, \qquad \xi_s \ge 0 && \forall s \in S \notag\\
& \vet x \in X \notag
\end{align}
All'ottimo $\eta^*$ è un VaR e il valore obiettivo è il CVaR: si ottengono
\emph{entrambe} le misure da un solo LP.
\end{modello}

\begin{spiegazione}
La variabile $\xi_s = (\ell_s - \eta)^+$ è l'eccesso di perdita oltre la soglia
candidata $\eta$: solo gli scenari nella coda contribuiscono (stesso trucco delle
parti positive del Newsvendor). Minimizzare su $\eta$ bilancia due spinte: alzare
$\eta$ riduce gli eccessi ma sposta in su la soglia; il punto di equilibrio è
esattamente il quantile $\alpha$. Se $X$ è un poliedro e $\ell_s$ è lineare, tutto il
modello è un LP.
\end{spiegazione}

\section{Esempio svolto a mano}

\begin{esempio}[Sei scenari, $\alpha = 0{,}80$]
Perdite equiprobabili $\{2, 4, 5, 7, 12, 20\}$.

\textbf{VaR.} La cumulata raggiunge $0{,}80$ in corrispondenza della quinta perdita
ordinata: $\var_{0,80} = 12$.

\textbf{CVaR con la formula.} La coda peggiore ha massa $0{,}20$; lo scenario da 20
pesa $1/6 = 0{,}167$, e la massa restante $0{,}033$ sta \emph{sul punto} 12:
\[
\cvar_{0,80} = 12 + \frac{1}{0{,}20} \cdot \frac{1}{6}\,(20 - 12)
= 12 + 6{,}67 = 18{,}67 .
\]

\textbf{Verifica con l'LP} (script): $\eta^* = 12$, $\xi = (0,0,0,0,0,8)$, valore
$18{,}6667$. La formula tratta correttamente la massa ``a cavallo'' del quantile
senza selezionare scenari con variabili binarie.

\textbf{Confronto.} Perdita media $8{,}33$; VaR $12$; CVaR $18{,}67$: tre numeri,
tre storie. La media rassicura, il VaR dà una soglia, il CVaR dice quanto fa male
quando va male.
\end{esempio}

\section{Caso di studio 1: portafoglio mean-CVaR contro Markowitz}

Stessi otto titoli del capitolo~\ref{cap:markowitz}, ma 220 scenari mensili generati
con \emph{code grasse} (fattore di mercato $t$ di Student, gradi di libertà 4):
il mondo in cui il CVaR mostra il suo valore. Perdita di scenario:
$\ell_s(\vet x) = -\sum_{i \in I} r_{is}\, x_i$, dove $r_{is}$ è il rendimento del titolo $i$ nello scenario $s$. Minimizziamo $\cvar_{0,90}$ con rendimento atteso
$\ge 8\%$ annuo.

\begin{lstlisting}[style=python]
eta = m.addVar(lb=-GRB.INFINITY, name="eta")   # soglia: libera!
xi  = m.addVars(S, name="xi")
m.addConstr(x.sum() == 1)
m.addConstr(gp.quicksum(mu[i]*x[i] for i in range(n)) >= r_min)
m.addConstrs((xi[s] >= -gp.quicksum(R[s, i]*x[i] for i in range(n))
              - eta for s in range(S)))
m.setObjective(eta + gp.quicksum(xi[s] for s in range(S))
                     / (S * (1 - alpha)), GRB.MINIMIZE)
\end{lstlisting}

\begin{lstlisting}[style=output]
            | perdita media | VaR90  | CVaR90   (perdite mensili)
  mean-CVaR |       -0,0067 | 0,0327 | 0,0489
  Markowitz |       -0,0067 | 0,0329 | 0,0531
Composizioni: mean-CVaR: ENE 9,9 FIN 1,9 TEC 10,5 SAN 8,6 UTL 46,5 MAT 22,5
              Markowitz: ENE 18,4 FIN 6,4 TEC 6,5 SAN 13,6 UTL 40,4 MAT 14,6
\end{lstlisting}

\begin{center}
\begin{tikzpicture}
\begin{axis}[labmezza, ybar, bar width=2.6pt,
  title={Distribuzione delle perdite (mean-CVaR)},
  xlabel={perdita mensile (\%)}, ylabel={scenari}, ymin=0]
\addplot [fill=teal!70, draw=none, forget plot] table [col sep=comma, x=centro, y=freq]
  {\figdat{cap13_istogramma}};
\addplot [verde, dashed, thick] coordinates {(-0.67, 0) (-0.67, 30)};
\addplot [arancio, dash dot, thick] coordinates {(3.27, 0) (3.27, 30)};
\addplot [rossomattone, thick] coordinates {(4.89, 0) (4.89, 30)};
\legend{media, VaR$_{90}$, CVaR$_{90}$}
\end{axis}
\end{tikzpicture}%
\begin{tikzpicture}
\begin{axis}[labmezza, title={Frontiera rendimento-CVaR},
  xlabel={CVaR$_{0{,}90}$ perdita mensile (\%)},
  ylabel={rendimento atteso annuo (\%)}]
\addplot+ [mark=*, mark size=1.4pt] table [col sep=comma, x=cvar, y=rend]
  {\figdat{cap13_frontiera}};
\end{axis}
\end{tikzpicture}
\captionof{figure}{A sinistra: istogramma delle perdite mensili del portafoglio
mean-CVaR sui 220 scenari, con media (verde), VaR$_{90}$ (arancione) e CVaR$_{90}$
(rosso): il CVaR è la media della coda a destra del VaR. A destra: frontiera
rendimento-CVaR ottenuta variando il rendimento minimo richiesto.}
\end{center}

\begin{insight}
A parità di rendimento richiesto e di perdita media, il portafoglio mean-CVaR taglia
la coda: CVaR$_{90}$ del 4,89\% contro 5,31\% di Markowitz ($-8\%$). Lo fa riducendo
i titoli ad alto beta (ENE dal 18\% al 10\%) a favore di UTL e --- sorpresa --- di una
quota di TEC che, pur mediocre in media, si muove in controtendenza negli scenari
peggiori di questo campione. La varianza penalizza simmetricamente sopra e sotto la
media; il CVaR guarda solo dove fa male.
\end{insight}

\section{Caso di studio 2: supply chain a due stadi}

Prima di osservare lo scenario si \emph{prenota} capacità dai fornitori; poi, visto
lo scenario, si decide il ricorso. F1 è economico ma fragile (nel 12\% degli scenari
la sua disponibilità crolla al 30\%); F2 è caro ma affidabile. Penalità di shortage:
40~\euro/unità; 400 scenari.

\begin{modello}[Due stadi con CVaR]
\begin{align*}
\min\;& (1-\lambda) \Bigl[\text{costo}_1(\vet x) + \textstyle\sum_{s \in S} \pi_s\,
        \text{costo}_2(\vet f_s, u_s)\Bigr] + \lambda\, \cvar_\alpha(\text{costo totale}) \\
\text{s.t.}\;\; & f_{as} \le \delta_{as}\, x_a \quad \forall a \in A,\ s \in S
   \qquad\text{(uso $\le$ capacità prenotata disponibile)} \\
& \textstyle\sum_{a \in A} f_{as} + u_s \ge d_s \quad \forall s \in S
   \qquad\text{(domanda servita o shortage)}
\end{align*}
Le prenotazioni $x_a$ sono \emph{di primo stadio} (uguali in tutti gli scenari);
flussi $f_{as}$ e shortage $u_s$ sono \emph{di ricorso}.
\end{modello}

\begin{spiegazione}
\textbf{Funzione obiettivo.} Combina il costo medio (prenotazioni più il valore
atteso di acquisti e penalità) con il CVaR del costo totale, pesati da $\lambda$.

\textbf{Primo vincolo.} Nello scenario $s$ si può usare il fornitore $a$ solo fino
alla capacità prenotata $x_a$ ridotta dalla disponibilità $\delta_{as} \in (0, 1]$
(il dato che modella il guasto).

\textbf{Secondo vincolo.} La domanda $d_s$ va coperta con i flussi o, in mancanza,
conteggiata come shortage $u_s$, che in obiettivo paga la penalità: il modello non è
mai inammissibile, ``compra'' l'inammissibilità al prezzo della penalità.
\end{spiegazione}

\begin{lstlisting}[style=output]
lambda = 0,0: prenoto F1 161,5  F2 134,2 | medio 1.724  CVaR90 2.902 | servizio 98,7%
lambda = 0,5: prenoto F1  92,7  F2 212,7 | medio 1.803  CVaR90 2.217 | servizio 99,9%
lambda = 1,0: prenoto F1  68,5  F2 236,9 | medio 1.906  CVaR90 2.139 | servizio 99,0%
\end{lstlisting}

\begin{insight}
Il decisore neutrale compra dal fornitore economico e accetta che nel 12\% degli
scenari lo shortage faccia esplodere i costi (CVaR 2.902). Al crescere di $\lambda$
la capacità migra verso il fornitore affidabile: $+79$~\euro{} di costo medio
comprano $-685$~\euro{} di CVaR. Questo è il \emph{costo della resilienza}, ed è
esattamente il numero che serve in una discussione su rishoring e dual sourcing.
\end{insight}

\section{Analisi di sensitività}

I parametri da variare, con l'effetto atteso:

\begin{center}\small
\begin{tabular}{lp{9.4cm}}
\toprule
$\alpha$ & più alto $\Rightarrow$ coda più estrema stimata da meno scenari:
           servono più scenari per stime stabili \\
$\lambda$ & controlla il compromesso prestazione media / protezione \\
$\pi_s$ & scenari ripesati o stress test (aggiungere scenari estremi) \\
$|S|$ & stabilità statistica: ripetere con campioni diversi, come nel
        capitolo~\ref{cap:newsvendor} \\
\bottomrule
\end{tabular}
\end{center}

\begin{attenzione}
Con $\alpha = 0{,}99$ e 220 scenari la coda contiene solo 2--3 scenari: il CVaR
stimato è quasi rumore. Regola pratica: servono almeno qualche decina di scenari
\emph{oltre} il quantile. Inoltre con distribuzioni discrete il VaR può non essere
unico ($\eta^*$ è \emph{un} VaR): confrontare sempre $\eta^*$ con il quantile
empirico.
\end{attenzione}

\section{Esercizi}

\begin{esercizio}[verifica]
Ricalcolare a mano VaR e CVaR dell'esempio a 6 scenari con $\alpha = 0{,}90$ e
verificare con l'LP. Attenzione: la coda ha massa $0{,}10$ e lo scenario da 20 pesa
$0{,}167 > 0{,}10$.
\end{esercizio}

\begin{esercizio}[subadditività]
Costruire un esempio a scenari con due attività in cui
$\var(L_1 + L_2) > \var(L_1) + \var(L_2)$ e verificare che il CVaR invece non viola
la subadditività.
\end{esercizio}

\begin{esercizio}[frontiera in $\alpha$]
Per il portafoglio, tracciare il CVaR ottimo per
$\alpha \in \{0{,}80;\, 0{,}90;\, 0{,}95;\, 0{,}99\}$ a parità di rendimento:
come cambia la composizione? Quando iniziano le instabilità numeriche?
\end{esercizio}

\begin{esercizio}[CVaR vincolato]
Riformulare il portafoglio come $\max \sum_{i \in I} \mu_i x_i$ s.t. $\cvar_{0,90} \le k$: verificare
che variando $k$ si percorre la stessa frontiera e leggere il duale del vincolo.
\end{esercizio}

\begin{esercizio}[stress test]
Aggiungere alla supply chain uno scenario estremo (domanda $+50\%$ e F1 al 10\%) con
probabilità 1\%: quanto cambia la prenotazione ottima per $\lambda = 0{,}5$?
E il CVaR?
\end{esercizio}
